\section*{Bab \ref{ch:b3:green-industrial} --- Kimia Hijau dan Proses Industri}

\begin{solution}{exo:b3:green-industrial:1}
Diels--Alder: 100\,\% (sebuah adisi). Esterifikasi: $88.0/(60.0 + 46.0) = 83$\,\%. Wittig:
$104.0/(106.0 + 276.0) = 27$\,\%, dengan trifenilfosfina oksida membawa pergi sebagian besar massa.
\end{solution}

\begin{solution}{exo:b3:green-industrial:2}
$\mathrm{PMI} = 120/8.0 = 15$; $E = \mathrm{PMI} - 1 = 14$.
\end{solution}

\begin{solution}{exo:b3:green-industrial:3}
Satu kali pencucian muatan standar pada suhu dan hasil kebersihan tertentu. Deterjen
didosis secara berbeda: serbuk pekat memerlukan lebih sedikit per pencucian, sehingga massa yang sama
tidak memberikan layanan yang sama.
\end{solution}

\begin{solution}{exo:b3:green-industrial:4}
2-Metiltetrahidrofuran atau etil asetat untuk diklorometana; heptana (dengan etil
asetat) untuk heksana; untuk DMF, etil asetat, 2-metiltetrahidrofuran, atau
dimetil karbonat, atau kopling di dalam air dengan \omterm{def:b3:colloids:surfactant}{surfaktan}, bergantung pada
pereaksinya.
\end{solution}

\begin{solution}{exo:b3:green-industrial:5}
1.000 mol asam; 0.800 mol ester (\qty{88.0}{g/mol}): rendemen 80\,\%. AE 83.0\,\%. SF $= 152/106 = 1.43$.
$\mathrm{RME} = 70.4/152 = 0.463$; $\mathrm{AE} \times \text{rendemen}/\mathrm{SF} = 0.830 \times 0.80/1.43 = 0.463$.
\end{solution}

\begin{solution}{exo:b3:green-industrial:6}
Klorohidrin: $44.0/(28.0 + 71.0 + 74.1) = 25$\,\%; oksidasi langsung: 100\,\%. Satu mol $\ce{CaCl2}$ per mol
etilena oksida: $\qty{1.0e6}{g}/\qty{44.0}{g/mol} \times \qty{111.1}{g/mol} = \qty{2.5}{t}$ per ton.
\end{solution}

\begin{solution}{exo:b3:green-industrial:7}
$3/8 \times \qty{5.88e4}{mol} \times \qty{44.0}{g/mol} = \qty{0.97}{t}$ $\ce{CO2}$ per ton amonia. Untuk hidrogen, $\ce{CH4 + 2H2O -> CO2 + 4H2}$:
$1/4$ mol per mol, $\qty{5.0e5}{mol}/4 \times \qty{44.0}{g/mol} = \qty{5.5}{t}$ $\ce{CO2}$ per ton hidrogen.
\end{solution}

\begin{solution}{exo:b3:green-industrial:8}
% ledger: dfg:H2O_l, dfh:H2O-l
Satu kilogram adalah \qty{500}{mol}. Dari $\Delta_rG^\circ = \qty{237.1}{kJ/mol}$: $\qty{1.19e5}{kJ} = \qty{32.9}{kWh}$. Dari $\Delta_rH^\circ = \qty{285.8}{kJ/mol}$: \qty{39.7}{kWh}
(selisihnya adalah kalor yang dapat dipasok oleh lingkungan).
\end{solution}

\begin{solution}{exo:b3:green-industrial:9}
$\qty{1.0e6}{g}/\qty{146.0}{g/mol} = \qty{6.85e3}{mol}$ $\ce{N2O}$, $\qty{0.30}{t}$; $\times 273 = \qty{82}{t}$ ekuivalen $\ce{CO2}$ per ton asam adipat.
\end{solution}

\begin{solution}{exo:b3:green-industrial:10}
RME sebuah tahap menghitung pereaksi berlebih sebagai terpakai; jika kelebihan itu diperoleh kembali
dan diumpankan kembali, hanya bagian yang benar-benar hilang yang terpakai, dan efisiensi massa
keseluruhannya lebih tinggi daripada hasil kali nilai-nilai tahapnya. Dalam \omterm{def:b3:green-industrial:e-factor}{faktor-E}, aliran
yang didaur ulang bukan limbah, asalkan daur ulang itu berada di dalam batas;
energi dan kehilangannya kemudian harus dihitung.
\end{solution}

\begin{solution}{exo:b3:green-industrial:11}
Sebelumnya: $20 \times 0.80 = \qty{16}{kg}$ pelarut, 10\,\% hilang: \qty{1.6}{kg} limbah. Sesudahnya: \qty{4.0}{kg}, \qty{0.40}{kg} hilang.
\omterm{def:b3:green-industrial:e-factor}{Faktor-E} lengkap turun 1.2 kg per kilogram produk (dan energi distilasi
turun tiga perempatnya).
\end{solution}

\begin{solution}{exo:b3:green-industrial:12}
Per \omterm{def:b3:green-industrial:lca}{satuan fungsional} dan dengan batas yang sama, LCA itu akan melaporkan setiap kategori
dampak secara terpisah (iklim, pemakaian lahan, pemakaian air\dots) dan bukan satu skor. Sebuah
keputusan memerlukan bobot yang diberikan kepada kategori-kategori itu, kelangkaan air
dan lahan setempat, ketidakpastian data, dan apakah bahan baku berbasis hayati itu
bersaing dengan pangan.
\end{solution}

\begin{solution}{pb:b3:green-industrial:1}
\textbf{1.} $\ce{C10H14 + C4H6O3 -> C12H16O + C2H4O2}$ (HF); $\ce{C12H16O + H2 -> C12H18O}$ (nikel Raney);
$\ce{C12H18O + CO -> C13H18O2}$ (paladium).

\textbf{2.} $\ce{C10H14 + C4H6O3 + H2 + CO -> C13H18O2 + C2H4O2}$.

\textbf{3.} $134.0 + 102.0 + 2.0 + 28.0 = \qty{266.0}{g/mol}$ pereaksi; ibuprofen \qty{206.0}{g/mol}: AE $= 206.0/266.0 = 77$\,\%.

\textbf{4.} Isobutilbenzena, anhidrida asetat, etil kloroasetat, natrium etoksida,
asam berair ($\ce{H3O+}$), hidroksilamina, dan air.

\textbf{5.} $134.0 + 102.0 + 122.5 + 68.0 + 19.0 + 33.0 + 2 \times 18.0 = \qty{514.5}{g/mol}$: AE $= 206.0/514.5 = 40$\,\%.

\textbf{6.} Asam asetat, natrium klorida, etanol, karbon dioksida, air, amonia (sebagai
garam amonium), dan garam aluminium dari aluminium klorida stoikiometrik.

\textbf{7.} $0.71 + 0.55 + 0.010 + 0.14 + 0.05 = \qty{1.46}{t}$.

\textbf{8.} PMI 1.46; $E = 0.46$.

\textbf{9.} $0.46 - 0.31 = 0.15$.

\textbf{10.} PMI $= 2.0 + 1.5 = 3.5$; $E = 2.5$.

\textbf{11.} Rendemen di bawah 100\,\%, pereaksi berlebih, kehilangan pelarut dan katalis, serta
bahan pengerjaan semuanya menambah limbah yang diabaikan oleh \omterm{def:b3:green-industrial:atom-economy}{ekonomi atom}.

\textbf{12.} Asas mencegah limbah, memaksimalkan \omterm{def:b3:green-industrial:atom-economy}{ekonomi atom}, dan memakai katalis
yang menggantikan pereaksi stoikiometrik, serta menghindari turunan dan tahap yang tidak perlu dan menghemat energi.

\textbf{13.} Aluminium klorida stoikiometrik, yang terhidrolisis menjadi limbah aluminium dalam
pengerjaan; hidrogen fluorida sekaligus menjadi katalis dan pelarut serta diperoleh kembali dan dipakai ulang.

\textbf{14.} Nikel Raney (katalis hidrogenasi heterogen).

\textbf{15.} Kompleks paladium fosfina: kimia \omterm{def:b3:homogeneous-catalysis:carbonylation}{karbonilasi} pada
\cref{ch:b3:homogeneous-catalysis}.

\textbf{16.} Anhidrida asetat memindahkan satu gugus asetil; separuh lain
molekulnya lepas sebagai asam asetat.

\textbf{17.} Hidrogen fluorida sangat beracun dan korosif, karbon monoksida beracun dan
mudah menyala: peralatan tertutup dari bahan yang tahan, detektor, dan operator
yang terlatih.

\textbf{18.} Setiap tahap menambah pemanasan, pendinginan, pemisahan, dan perolehan kembali pelarut; tiga
tahap memerlukan lebih sedikit dari semua itu.

\textbf{19.} $2.5 \times 3000 = \qty{7500}{t}$ per tahun.

\textbf{20.} $0.15 \times 3000 = \qty{450}{t}$ per tahun.

\textbf{21.} Sekitar \qty{7000}{t} limbah yang dihindari setiap tahun.

\textbf{22.} Tidak: jenis limbahnya penting (satu ton air garam bukan satu ton
senyawa beracun yang persisten), demikian pula energi, emisi, dan bahaya
pereaksi-pereaksinya.

\textbf{23.} Dampak pembuatan pereaksi dan energi di hulu, dampak emisi ke
udara dan air, serta dampak pemakaian dan pembuangan produk, per \omterm{def:b3:green-industrial:lca}{satuan fungsional}, dalam
beberapa kategori dampak.

\textbf{24.} Rute katalitik tiga tahap memiliki \omterm{def:b3:green-industrial:atom-economy}{ekonomi atom} sekitar 77\,\% (40\,\% untuk
rute enam tahap).
\end{solution}
